% Authoritative LaTeX chapter. Edit this file for future revisions.
% Initially migrated 2026-09-11; no Markdown import occurs during PDF builds.
\chapter{Python 사용과 배치 자료 생성}
\label{chap:06_python_batch}
매뉴얼 목차 · \hyperref[chap:01_concepts_geometry]{RAA와 관측 기하} · \hyperref[chap:04_examples]{C 실행 예시}

이 부록은 \texttt{code/\allowbreak{}OCRT\_\allowbreak{}Python/\allowbreak{}}의 공개 실행 스크립트와 \texttt{code/\allowbreak{}campaign\_\allowbreak{}runner/\allowbreak{}}의 C 실행 관리 도구를 설명한다. 옵션·기본값·출력은 2026-09-11에 확인한 소스 기준이다. \textbf{Python 계산기는 현재 C v1.11.1의 모든 기능을 제공하지 않는다.} C용 옵션을 Python 명령에 그대로 옮기거나, 과거 Python README의 일치도 수치를 현재 자료 전체에 적용하지 않는다.

\section{먼저 알아둘 지원 범위}
\label{sec:auto:06_python_batch:1}

{\TableFont
\begin{longtable}{@{}>{\raggedright\arraybackslash}p{\dimexpr 0.3000\linewidth-5.33pt\relax}>{\raggedright\arraybackslash}p{\dimexpr 0.2300\linewidth-5.33pt\relax}>{\raggedright\arraybackslash}p{\dimexpr 0.4700\linewidth-5.33pt\relax}@{}}
\toprule
\textbf{목적} & \textbf{사용할 진입점} & \textbf{특성} \\
\midrule
\endfirsthead
\toprule
\textbf{목적} & \textbf{사용할 진입점} & \textbf{특성} \\
\midrule
\endhead
\bottomrule
\endfoot
한 조건의 R1/R2/R3 확인 & \texttt{ocrt\_\allowbreak{}solve.\allowbreak{}py single} & 계산할 반사도 성분 선택, 압력 지정, 수렴 상태 표시 \\
여러 조건을 순차 계산 & \texttt{ocrt\_\allowbreak{}solve.\allowbreak{}py grid} & 행별 단일 계산, 일부 행·밴드 선택, 명시적 재개 \\
여러 조건을 NumPy/CuPy 배치 계산 & \texttt{produce\_\allowbreak{}grid.\allowbreak{}py} & 각 조건마다 R1/R2/R3 모두 계산, 자동 재개, 상세 파생 출력 \\
같은 물리 조건에서 VZA \ensuremath{\times} RAA 각도 격자 & \texttt{ocrt\_\allowbreak{}fullgrid.\allowbreak{}py} & 대기·수중장을 공유하는 각도 LUT, I/Q/U 출력 \\
CPU/GPU 결과 대조 & \texttt{gpu\_\allowbreak{}gate.\allowbreak{}py} & 네 가지 I 성분 제품 대조 \\
C 계산을 여러 프로세스로 관리 & \texttt{campaign\_\allowbreak{}runner/\allowbreak{}runner/\allowbreak{}run\_\allowbreak{}campaign.\allowbreak{}py} & Python이 C 실행파일을 호출하는 별도 도구 \\
\end{longtable}
}

현재 버전에서 특히 다음을 구분한다.

\begin{itemize}
\item \textbf{구면 보정:} 위 Python 계산 CLI에는 \texttt{-\allowbreak{}-\allowbreak{}pssa}나 IPSS 옵션이 없다. 최신 C의 IPSS 적용 결과가 필요한 LUT·자료 생산은 C 경로를 사용한다. Python을 C의 기본값까지 완전히 같은 대체 구현으로 취급하지 않는다.
\item \textbf{해면:} Python 결합 단일/각도 LUT 경로는 \texttt{wind <=\allowbreak{} 0}일 때 평탄 해면 결합이 이식되지 않았다는 \texttt{NotImplementedEr\allowbreak{}ror}를 낸다. 예시는 양의 풍속을 사용한다. 배치 경로에서 숫자 0을 받는 것과 평탄 해면 구현을 지원하는 것은 같은 의미가 아니다.
\item \textbf{Chl:} 종 옵션을 생략하면 기본 흡수표를 쓰고 식물플랑크톤 자체 \texttt{b =\allowbreak{} bb =\allowbreak{} 0}으로 처리한다. Chl 연동 쇄설물 산란은 별도로 남는다. \texttt{Chl > 0}에서 \texttt{-\allowbreak{}-\allowbreak{}phyto}/\texttt{-\allowbreak{}-\allowbreak{}phyto-\allowbreak{}group}를 명시하면 기본 종 이름이라도 오류다. \texttt{OCRT\_\allowbreak{}ADVANCED=\allowbreak{}1}로 해제되지 않는다.
\item \textbf{설치 자료 누락:} 현재 생성자는 사용하지 않는 EAP 종별 Mie도 읽는다. 기본 경로는 \texttt{data/\allowbreak{}water\_\allowbreak{}iop/\allowbreak{}eap/\allowbreak{}EAP\_\allowbreak{}02\_\allowbreak{}Diatoms\_\allowbreak{}centric\_\allowbreak{}D6.\allowbreak{}mie}이다. 이 파일은 현재 \texttt{data-\allowbreak{}v1}의 181개 파일 매니페스트에 포함되지 않는다. 매뉴얼 작성 시 로컬에서도 이 파일 누락으로 성분 모델 초기화가 실패했다. 따라서 \textbf{현재 저장소와 공개 181개 자료를 설치하는 것만으로 Python R1/배치/각도 LUT 실행 준비가 끝난다고 볼 수 없다.} 배포자에게 일치하는 EAP 보조 자료를 받아 경로와 이력을 확인하거나, C 경로를 사용한다. 다른 Mie를 해당 이름으로 바꿔 넣지 않는다.
\item \textbf{배치·단일의 수치 경로:} \texttt{produce\_\allowbreak{}grid.\allowbreak{}py}의 R1 대기 패스는 Greek 계수 커널, R3는 value 커널을 사용한다. 단일 CLI의 R1/R3는 value 커널을 사용한다. 배치 R2는 별도 층 수·산란차수 기본값도 사용한다. 그러므로 CPU/GPU 일치, 배치/단일 일치, 현재 C와의 일치를 각각 확인해야 한다.
\item \textbf{파장:} 기본 파장 계약은 330--1100 nm지만 실제 사용되는 흡수·위상 표가 해당 파장을 덮어야 한다. 오래된 소스 주석에 남은 쇄설물 350--850 nm 범위와 실제 설치된 표가 다를 수 있다. 현재 로컬 FR631 쇄설물 위상표는 읽기 검사에서 330--1100 nm였다. 자료 버전과 SHA-256을 기록한다.
\end{itemize}

\texttt{produce\_\allowbreak{}grid.\allowbreak{}py}의 파일 머리에 남아 있는 “기체흡수 미적용”, “R2/R3에 value 커널 미적용”, \texttt{-\allowbreak{}-\allowbreak{}strict-\allowbreak{}check} 안내는 현 실행 코드와 맞지 않는다. 현재 기체흡수는 배치 대기에 전달되고, R3는 value 커널을 쓰며, \texttt{-\allowbreak{}-\allowbreak{}strict-\allowbreak{}check}라는 CLI 옵션은 없다.

\section{설치와 준비 검사}
\label{sec:auto:06_python_batch:2}
아래에서 명령은 저장소의 \texttt{code/\allowbreak{}OCRT\_\allowbreak{}Python} 폴더에서 실행한다. \texttt{ocrt\_\allowbreak{}py}는 이 폴더에서 가져오는 모듈이며, 현재 별도의 배포용 \texttt{pyproject.\allowbreak{}toml}/\texttt{setup.\allowbreak{}py}가 없다.

CPU 준비 예시:


\begin{ManualCode}
cd <OCRT 저장소 경로>\code\OCRT_Python
python -m venv .venv
.\.venv\Scripts\python -m pip install "numpy>=2" scipy
.\.venv\Scripts\python -c "import numpy, scipy; print(numpy.__version__, scipy.__version__)"
\end{ManualCode}

Linux/WSL에서는 \texttt{python3 -\allowbreak{}m venv .\allowbreak{}venv}, \texttt{.\allowbreak{}venv/\allowbreak{}bin/\allowbreak{}python -\allowbreak{}m pip install 'n\allowbreak{}umpy>=\allowbreak{}2' scipy}처럼 실행한다. 이후 예시의 \texttt{python}은 준비한 가상환경의 Python으로 바꿔 사용한다. Python 3.10 이상을 기준으로 환경을 구성하면 이 부록의 계산기와 캠페인 도구를 함께 쓰기 편하다. NumPy 2 이상을 적은 이유는 현재 코드가 \texttt{numpy.\allowbreak{}trapezoid}를 사용하기 때문이고, SciPy는 CPU 배치의 \texttt{erfc} 계산에 필요하다. 예전 README의 “NumPy만 설치” 안내는 이 배치 경로에는 충분하지 않다.

Mie 자료는 저장소 루트의 \texttt{scripts/\allowbreak{}fetch\_\allowbreak{}data.\allowbreak{}ps1} 또는 \texttt{scripts/\allowbreak{}fetch\_\allowbreak{}data.\allowbreak{}sh}로 설치한다. 파일은 C \texttt{inputs/\allowbreak{}}와 Python \texttt{data/\allowbreak{}}에 같은 상대경로로 놓인다. \textbf{이 설치와 앞 절의 EAP 누락 확인은 별개다.} 에어로졸 AOD가 0이어도 Python 고수준 CLI는 \texttt{.\allowbreak{}mie}를 읽는다. C Rayleigh 단독 실행의 “Mie 없이 실행” 조건을 Python 예시에 적용하지 않는다.

준비 검사:


\begin{ManualCode}
python produce_grid.py --help
python ocrt_solve.py single --help
python ocrt_fullgrid.py --help
python -c "from pathlib import Path; print(Path('data/water_iop/eap/EAP_02_Diatoms_centric_D6.mie').is_file())"
python -c "from ocrt_py.constituent import OCRTConstituentModel; OCRTConstituentModel('data'); print('constituent data loaded')"
\end{ManualCode}

마지막 검사가 성공해야 R1 예제를 진행할 수 있다. \texttt{-\allowbreak{}-\allowbreak{}help} 성공은 자료 준비나 복사전달 계산 성공을 뜻하지 않는다.

GPU는 호환되는 NVIDIA CUDA 환경과 그 환경에 맞는 CuPy가 추가로 필요하다. 예를 들어 CUDA 12 계열 환경에서는 \texttt{python -\allowbreak{}m pip install cu\allowbreak{}py-\allowbreak{}cuda12x}를 사용한다. 실제 드라이버/CUDA/Python 조합은 \href{https://docs.cupy.dev/en/stable/install.html}{CuPy 공식 설치 지침}과 맞춘다. GPU 선택은 \textbf{Python 프로세스를 시작하기 전에} 설정한다.


\begin{ManualCode}
$env:OCRT_PY_GPU = "1"
python produce_grid.py --grid grid.csv --out result_gpu.csv --bands 443,555 --chunk 4
Remove-Item Env:OCRT_PY_GPU
\end{ManualCode}

Linux/WSL에서는 \texttt{OCRT\_\allowbreak{}PY\_\allowbreak{}GPU=\allowbreak{}1 python produce\allowbreak{}\_\allowbreak{}grid.\allowbreak{}py .\allowbreak{}.\allowbreak{}.\allowbreak{}}이다. 값이 정확히 \texttt{1}일 때 CuPy를 가져온다. CuPy가 없으면 오류이며 자동 CPU 전환은 하지 않는다. \texttt{ocrt\_\allowbreak{}solve.\allowbreak{}py}와 각도 LUT의 모든 단계가 GPU화되었다는 의미로 해석하지 않는다.

\Needspace{7\baselineskip}
\section{공통 입력 규약: 단위와 RAA}
\label{sec:auto:06_python_batch:3}

{\TableFont
\begin{longtable}{@{}>{\raggedright\arraybackslash}p{\dimexpr 0.3400\linewidth-4.00pt\relax}>{\raggedright\arraybackslash}p{\dimexpr 0.6600\linewidth-4.00pt\relax}@{}}
\toprule
\textbf{입력} & \textbf{의미·단위} \\
\midrule
\endfirsthead
\toprule
\textbf{입력} & \textbf{의미·단위} \\
\midrule
\endhead
\bottomrule
\endfoot
\texttt{sza} & 태양천정각, 도 \\
\texttt{vza} & 공기 쪽 관측천정각, 도 \\
\texttt{raa} & OCRT 상대방위각, 도. \textbf{180\textdegree{}가 경면/선글린트 쪽, 0\textdegree{}가 후방산란 쪽} \\
\texttt{wind} & 풍속, m/s \\
\texttt{aerosol} 또는 \texttt{aer} & \texttt{data/\allowbreak{}} 기준 Mie 경로 또는 모델 이름 \\
\texttt{aod865} & 865 nm 기준 에어로졸 광학두께, 무차원 \\
\texttt{chl} & Chl 농도, mg/m\textsuperscript{3} \\
\texttt{tsm} & 부유물 농도, g/m\textsuperscript{3}. 수치상 mg/L와 같음 \\
\texttt{acdom440} & 440 nm CDOM 흡수계수, m\textsuperscript{-}\textsuperscript{1} \\
\end{longtable}
}

\textbf{다른 프로그램의 RAA를 그대로 복사하기 전에 정의를 확인한다.} 비교 자료가 경면 방향을 \texttt{0\textdegree{}}로 놓는 대칭 범위 \texttt{0--180\textdegree{}} 규약이면 OCRT 입력은 \texttt{180\textdegree{} - RAA\_\allowbreak{}외부}이다. 예를 들어 외부 \texttt{0/\allowbreak{}60/\allowbreak{}180\textdegree{}}는 OCRT \texttt{180/\allowbreak{}120/\allowbreak{}0\textdegree{}}가 된다. 이 변환은 두 규약의 방향 정의를 확인한 경우에만 적용한다. 위성 태양·센서 방위각에서 부호 있는 \texttt{0--360\textdegree{}} RAA를 만들 때는 기준 벡터와 회전 방향까지 기록해야 한다. \textbf{Q/U를 다룰 때 RAA를 임의로 \texttt{0--180\textdegree{}}로 접으면 U의 부호 정보가 사라질 수 있다.} 예시에는 \texttt{90\textdegree{}}와 \texttt{270\textdegree{}}를 함께 넣어 방위 대칭을 확인할 수 있게 했다.

일반적인 상향 관측 예시는 SZA/VZA를 \texttt{0 <=\allowbreak{} 각도 < 90\textdegree{}} 범위에 둔다. \texttt{make\_\allowbreak{}grid.\allowbreak{}py}의 65\textdegree{}/55\textdegree{} 상한은 해당 샘플 설계의 범위이며 전체 OCRT의 검증 범위를 선언한 값이 아니다. 정확한 천저 \texttt{VZA=\allowbreak{}0\textdegree{}}는 일부 선택적 polarized value 수중 커널의 출력 특이점이 알려져 있다. 그 커널을 연구 목적으로 켰다면 \texttt{0.\allowbreak{}001\textdegree{}} 등 명시적 비천저 조건으로 비교하고 대체 각도를 기록한다.

\subsection{배치 CSV}
\label{sec:auto:06_python_batch:3.1}
\texttt{produce\_\allowbreak{}grid.\allowbreak{}py}의 필수 헤더는 다음 10개다. 헤더 대소문자는 구분하지 않지만 이름 안의 공백·철자까지 자동 정리하지 않는다. 숫자는 소수점 \texttt{.\allowbreak{}}을 쓰고 단위 문자열은 넣지 않는다. \texttt{case\_\allowbreak{}id}는 파일 내에서 고유해야 한다.


\begin{ManualCode}
case_id,sza,vza,raa,wind,aerosol,aod865,chl,tsm,acdom440
open_ocean,30,20,90,5,C50,0.05,0.2,0.1,0.01
coastal,40,35,90,5,M80C,0.10,1.0,3.0,0.04
\end{ManualCode}

\texttt{produce\_\allowbreak{}grid.\allowbreak{}py}는 이름 끝에 \texttt{.\allowbreak{}mie}를 항상 붙이므로 \texttt{C50.\allowbreak{}mie} 대신 \textbf{\texttt{C50}}을 쓴다. 하위 폴더 모델은 \texttt{aerosol\_\allowbreak{}ahmad2010\_\allowbreak{}paper\_\allowbreak{}mie/\allowbreak{}A2010ver\_\allowbreak{}r80f50v01}처럼 쓴다. 단일/각도 CLI의 \texttt{-\allowbreak{}-\allowbreak{}aer}는 \texttt{.\allowbreak{}mie} 확장자가 있어도 된다. 실제 설치된 모델명을 사용하며, 이름이 비슷한 AccuRT 계열과 Ahmad 논문 기반 계열을 혼용하지 않는다.

선택적 CSV 열은 비어 있으면 기본값을 사용한다.


{\TableFont
\begin{longtable}{@{}>{\raggedright\arraybackslash}p{\dimexpr 0.3000\linewidth-5.33pt\relax}>{\raggedright\arraybackslash}p{\dimexpr 0.2300\linewidth-5.33pt\relax}>{\raggedright\arraybackslash}p{\dimexpr 0.4700\linewidth-5.33pt\relax}@{}}
\toprule
\textbf{대표 헤더} & \textbf{허용되는 다른 이름 (\texttt{produce\_\allowbreak{}grid.\allowbreak{}py})} & \textbf{기본값·효과} \\
\midrule
\endfirsthead
\toprule
\textbf{대표 헤더} & \textbf{허용되는 다른 이름 (\texttt{produce\_\allowbreak{}grid.\allowbreak{}py})} & \textbf{기본값·효과} \\
\midrule
\endhead
\bottomrule
\endfoot
\texttt{ocrt\_\allowbreak{}adom\_\allowbreak{}slope} & \texttt{adom\_\allowbreak{}slope}, \texttt{cdom\_\allowbreak{}slope} & 0.014 nm\textsuperscript{-}\textsuperscript{1}. \texttt{aCDOM(\ensuremath{\lambda})=\allowbreak{}aCDOM440·exp[-S(\allowbreak{}\ensuremath{\lambda}-440)]} \\
\texttt{ocrt\_\allowbreak{}detritus\_\allowbreak{}a440} & \texttt{detritus\_\allowbreak{}a440}, \texttt{det\_\allowbreak{}a440} & 0 m\textsuperscript{-}\textsuperscript{1}. Chl 연동 쇄설물의 440 nm 흡수 크기 \\
\texttt{ocrt\_\allowbreak{}detritus\_\allowbreak{}slope} & \texttt{detritus\_\allowbreak{}slope}, \texttt{det\_\allowbreak{}slope} & 0.0109 nm\textsuperscript{-}\textsuperscript{1}. 쇄설물 흡수의 지수 기울기 \\
\end{longtable}
}

\texttt{ocrt\_\allowbreak{}solve.\allowbreak{}py grid}는 같은 대표 헤더를 읽으며 추가 별칭도 지원한다. 예를 들어 ID는 \texttt{row\_\allowbreak{}id/\allowbreak{}id/\allowbreak{}case/\allowbreak{}index/\allowbreak{}row}, 에어로졸은 \texttt{aer\_\allowbreak{}model/\allowbreak{}aerosol\_\allowbreak{}model/\allowbreak{}mie/\allowbreak{}model/\allowbreak{}aer}, 각도는 \texttt{sza\_\allowbreak{}deg/\allowbreak{}vza\_\allowbreak{}deg/\allowbreak{}raa\_\allowbreak{}deg}, TSM은 \texttt{tsm\_\allowbreak{}g\_\allowbreak{}m3/\allowbreak{}tsm\_\allowbreak{}gm3}, Chl은 \texttt{chl\_\allowbreak{}mg\_\allowbreak{}m3/\allowbreak{}chl\_\allowbreak{}mgm3/\allowbreak{}chla}, CDOM은 \texttt{adom440/\allowbreak{}a\_\allowbreak{}cdom\_\allowbreak{}440/\allowbreak{}acdom440\_\allowbreak{}m\_\allowbreak{}inv/\allowbreak{}acdom\_\allowbreak{}440/\allowbreak{}cdom440/\allowbreak{}a\_\allowbreak{}cdom440}, AOD는 \texttt{aod\_\allowbreak{}865}, 풍속은 \texttt{wind\_\allowbreak{}ms/\allowbreak{}wind\_\allowbreak{}speed/\allowbreak{}wind\_\allowbreak{}speed\_\allowbreak{}ms/\allowbreak{}wind\_\allowbreak{}m\_\allowbreak{}s}도 가능하다. 헤더 정규화는 영문·숫자만 남겨 비교한다. ID 생략 시 행 순번을 사용한다. 여러 도구에서 함께 사용할 CSV는 위의 대표 헤더를 유지하는 편이 안전하다.

\section{\texttt{produce\_\allowbreak{}grid.\allowbreak{}py}: 전체 옵션}
\label{sec:auto:06_python_batch:4}

{\TableFont
\begin{longtable}{@{}>{\raggedright\arraybackslash}p{\dimexpr 0.3000\linewidth-5.33pt\relax}>{\raggedright\arraybackslash}p{\dimexpr 0.2300\linewidth-5.33pt\relax}>{\raggedright\arraybackslash}p{\dimexpr 0.4700\linewidth-5.33pt\relax}@{}}
\toprule
\textbf{옵션} & \textbf{기본값} & \textbf{설명} \\
\midrule
\endfirsthead
\toprule
\textbf{옵션} & \textbf{기본값} & \textbf{설명} \\
\midrule
\endhead
\bottomrule
\endfoot
\texttt{-\allowbreak{}h}, \texttt{-\allowbreak{}-\allowbreak{}help} & --- & 도움말 출력 후 종료 \\
\texttt{-\allowbreak{}-\allowbreak{}grid PATH} & 필수 & 입력 조건 CSV \\
\texttt{-\allowbreak{}-\allowbreak{}out PATH} & 필수 & 결과 CSV. 존재하면 자동으로 이어 쓰며 완료 키를 건너뜀 \\
\texttt{-\allowbreak{}-\allowbreak{}data DIR} & 스크립트 옆 \texttt{data/\allowbreak{}} & 에어로졸·해수 IOP·AFGL·기체 단면적 자료 루트 \\
\texttt{-\allowbreak{}-\allowbreak{}phyto NAME} & 생략 & 기본 Chl 흡수 경로. Chl > 0 행이 있으면 명시적 종 지정 금지 \\
\texttt{-\allowbreak{}-\allowbreak{}mineral NAME} & \texttt{red\_\allowbreak{}clay} & 광물 종류. \texttt{red\_\allowbreak{}clay}, \texttt{brown\_\allowbreak{}earth}, \texttt{yellow\_\allowbreak{}clay}, \texttt{calcareous\_\allowbreak{}sand} \\
\texttt{-\allowbreak{}-\allowbreak{}chunk N} & 64 & 한 배치의 조건 행 수. 양의 정수. GPU 메모리·케이스 복잡도에 맞춰 조절 \\
\texttt{-\allowbreak{}-\allowbreak{}n-\allowbreak{}mu-\allowbreak{}water N} & 24 & 수중과 배치 대기 Gauss 적분 절점 수에 전달됨 \\
\texttt{-\allowbreak{}-\allowbreak{}nt-\allowbreak{}atm N} & 400 & R1/R3 대기 SOS 층 수 \\
\texttt{-\allowbreak{}-\allowbreak{}m-\allowbreak{}max N} & 16 & R1/R3 대기 Fourier 모드 상한 \\
\texttt{-\allowbreak{}-\allowbreak{}max-\allowbreak{}it-\allowbreak{}water N} & 500 & 수중 산란차수 상한. 상한 도달과 수렴은 별개 \\
\texttt{-\allowbreak{}-\allowbreak{}bands LIST} & 아래 12밴드 & 쉼표로 구분한 정수 nm. 예: \texttt{443,\allowbreak{}555,\allowbreak{}865} \\
\texttt{-\allowbreak{}-\allowbreak{}no-\allowbreak{}gas} & 미지정, 흡수 켜짐 & 기체흡수 끄기 \\
\end{longtable}
}

기본 밴드는 \texttt{380,\allowbreak{}412,\allowbreak{}443,\allowbreak{}490,\allowbreak{}510,\allowbreak{}555,\allowbreak{}620,\allowbreak{}660,\allowbreak{}680,\allowbreak{}709,\allowbreak{}745,\allowbreak{}865}이다. 이 목록을 특정 센서의 최종 응답함수 적분 결과와 혼동하지 않는다. 계산은 지정한 파장점에서 이루어진다.

CLI로 바꾸지 못하는 주요 조건도 결과와 함께 기록해야 한다: 압력 \texttt{1013.\allowbreak{}25 hPa}, 물 굴절률 \texttt{1.\allowbreak{}34}, R1/R3 대기 산란차수 상한 \texttt{100}, 대기/수중 허용오차 \texttt{1e-\allowbreak{}7}, 에어로졸 전개 상한 \texttt{80}, 수중 입자 전개 상한 \texttt{200}이다. R2는 별도 함수 기본값인 \textbf{40층, Fourier 상한 2, 최대 20차}를 사용하며 \texttt{-\allowbreak{}-\allowbreak{}nt-\allowbreak{}atm}, \texttt{-\allowbreak{}-\allowbreak{}m-\allowbreak{}max}를 바꿔도 함께 바뀌지 않는다. 기체흡수 모델은 기본 \texttt{usstd76}이며 이 CLI에는 대기 프로파일·개별 기체량 지정 옵션이 없다.

\subsection{재개 규칙}
\label{sec:auto:06_python_batch:4.1}
기존 CSV의 헤더가 현재 \texttt{OUT\_\allowbreak{}FIELDS}와 정확히 같아야 한다. 맞으면 \texttt{(case\_\allowbreak{}id,\allowbreak{} band\_\allowbreak{}nm)}만으로 완료 여부를 판단한다. 입력 값·자료 파일·광물 종류·기체 설정·수치 옵션이 달라졌는지는 비교하지 않으며, 기록된 NaN이나 미수렴 결과도 자동 재계산하지 않는다. 한 CSV를 동시에 여러 프로세스에서 쓰지 않는다.

\textbf{입력이나 계산 설정을 바꾸면 새 출력 파일을 사용한다.} 중단 후 재개는 동일한 입력·설정·자료일 때만 한다. 출력 부모 폴더는 먼저 만든다. 출력은 밴드 순서 \ensuremath{\rightarrow} 청크 \ensuremath{\rightarrow} 입력 행 순서이며 행 ID 순서로 정렬된다고 가정하지 않는다. 입력 CSV의 중복 ID나 중복 밴드는 실행 전에 제거한다.

\section{\texttt{ocrt\_\allowbreak{}solve.\allowbreak{}py}: 단일·순차 격자 전체 옵션}
\label{sec:auto:06_python_batch:5}
사용 형식은 \texttt{python ocrt\_\allowbreak{}solve.\allowbreak{}py single .\allowbreak{}.\allowbreak{}.\allowbreak{}} 또는 \texttt{python ocrt\_\allowbreak{}solve.\allowbreak{}py grid .\allowbreak{}.\allowbreak{}.\allowbreak{}}이다. 공통 옵션은 서브명령 뒤에 둔다.


{\TableFont
\begin{longtable}{@{}>{\raggedright\arraybackslash}p{\dimexpr 0.3000\linewidth-5.33pt\relax}>{\raggedright\arraybackslash}p{\dimexpr 0.2300\linewidth-5.33pt\relax}>{\raggedright\arraybackslash}p{\dimexpr 0.4700\linewidth-5.33pt\relax}@{}}
\toprule
\textbf{공통 옵션} & \textbf{기본값} & \textbf{설명} \\
\midrule
\endfirsthead
\toprule
\textbf{공통 옵션} & \textbf{기본값} & \textbf{설명} \\
\midrule
\endhead
\bottomrule
\endfoot
\texttt{-\allowbreak{}h}, \texttt{-\allowbreak{}-\allowbreak{}help} & --- & 해당 서브명령 도움말 \\
\texttt{-\allowbreak{}-\allowbreak{}which R1 R2 R3} & 모두 & 실행할 성분을 공백으로 구분. \texttt{-\allowbreak{}-\allowbreak{}which R2}로 Rayleigh 기준만 계산 가능 \\
\texttt{-\allowbreak{}-\allowbreak{}phyto-\allowbreak{}group NAME} & 생략 & 종 선택 제한은 \ref{sec:auto:06_python_batch:1}절과 같음 \\
\texttt{-\allowbreak{}-\allowbreak{}tsm-\allowbreak{}species NAME} & \texttt{red\_\allowbreak{}clay} & 광물 네 종류 중 선택 \\
\texttt{-\allowbreak{}-\allowbreak{}n-\allowbreak{}mu-\allowbreak{}water N} & 48 & 결합 R1 수중 Gauss 절점 수. R2/R3 대기 절점 수는 내부 24로 고정 \\
\texttt{-\allowbreak{}-\allowbreak{}m-\allowbreak{}max N} & 16 & R1/R3 대기 Fourier 상한. R2는 2로 고정 \\
\texttt{-\allowbreak{}-\allowbreak{}n-\allowbreak{}layers N} & 400 & R1/R2/R3 대기 층 수 \\
\texttt{-\allowbreak{}-\allowbreak{}max-\allowbreak{}it-\allowbreak{}water N} & 500 & R1 수중 산란차수 상한 \\
\texttt{-\allowbreak{}-\allowbreak{}pressure HPA} & 1013.25 & Rayleigh 계산에 전달할 압력, hPa \\
\texttt{-\allowbreak{}-\allowbreak{}no-\allowbreak{}gas} & 미지정 & 기체흡수 끄기. 압력 0만 주어도 기체가 꺼지는 것은 아님 \\
\texttt{-\allowbreak{}-\allowbreak{}adom-\allowbreak{}slope S}, \texttt{-\allowbreak{}-\allowbreak{}cdom-\allowbreak{}slope S} & 0.014 & CDOM 지수 흡수 기울기, nm\textsuperscript{-}\textsuperscript{1}. 같은 옵션의 별칭 \\
\texttt{-\allowbreak{}-\allowbreak{}detritus-\allowbreak{}a440 A} & 0 & Chl 연동 쇄설물 흡수 크기, m\textsuperscript{-}\textsuperscript{1} \\
\texttt{-\allowbreak{}-\allowbreak{}detritus-\allowbreak{}slope S} & 0.0109 & 쇄설물 흡수 기울기, nm\textsuperscript{-}\textsuperscript{1} \\
\end{longtable}
}

\texttt{single} 전용 옵션:


{\TableFont
\begin{longtable}{@{}>{\raggedright\arraybackslash}p{\dimexpr 0.3000\linewidth-5.33pt\relax}>{\raggedright\arraybackslash}p{\dimexpr 0.2300\linewidth-5.33pt\relax}>{\raggedright\arraybackslash}p{\dimexpr 0.4700\linewidth-5.33pt\relax}@{}}
\toprule
\textbf{옵션} & \textbf{기본값} & \textbf{설명} \\
\midrule
\endfirsthead
\toprule
\textbf{옵션} & \textbf{기본값} & \textbf{설명} \\
\midrule
\endhead
\bottomrule
\endfoot
\texttt{-\allowbreak{}-\allowbreak{}wl NM} & 필수 & 파장, nm. 실수 가능 \\
\texttt{-\allowbreak{}-\allowbreak{}sza DEG}, \texttt{-\allowbreak{}-\allowbreak{}vza DEG}, \texttt{-\allowbreak{}-\allowbreak{}raa DEG} & 각각 필수 & 관측 기하, \ref{sec:auto:06_python_batch:3}절의 OCRT 규약 \\
\texttt{-\allowbreak{}-\allowbreak{}wind MPS} & 필수 & 풍속, 결합 R1은 양수 필요 \\
\texttt{-\allowbreak{}-\allowbreak{}chl VALUE} & 0 & Chl 농도, mg/m\textsuperscript{3} \\
\texttt{-\allowbreak{}-\allowbreak{}tsm VALUE} & 0 & TSM 농도, g/m\textsuperscript{3} \\
\texttt{-\allowbreak{}-\allowbreak{}acdom440 VALUE} & 0 & CDOM 흡수계수, m\textsuperscript{-}\textsuperscript{1} \\
\texttt{-\allowbreak{}-\allowbreak{}aod865 VALUE} & 0 & 865 nm AOD \\
\texttt{-\allowbreak{}-\allowbreak{}aer NAME} & \texttt{C50} & \texttt{data/\allowbreak{}} 기준 모델명/상대경로. \texttt{.\allowbreak{}mie} 확장자 선택 가능 \\
\end{longtable}
}

\texttt{grid} 전용 옵션:


{\TableFont
\begin{longtable}{@{}>{\raggedright\arraybackslash}p{\dimexpr 0.3000\linewidth-5.33pt\relax}>{\raggedright\arraybackslash}p{\dimexpr 0.2300\linewidth-5.33pt\relax}>{\raggedright\arraybackslash}p{\dimexpr 0.4700\linewidth-5.33pt\relax}@{}}
\toprule
\textbf{옵션} & \textbf{기본값} & \textbf{설명} \\
\midrule
\endfirsthead
\toprule
\textbf{옵션} & \textbf{기본값} & \textbf{설명} \\
\midrule
\endhead
\bottomrule
\endfoot
\texttt{-\allowbreak{}-\allowbreak{}grid PATH}, \texttt{-\allowbreak{}-\allowbreak{}out PATH} & 각각 필수 & 입력 조건 CSV와 결과 CSV \\
\texttt{-\allowbreak{}-\allowbreak{}bands NM .\allowbreak{}.\allowbreak{}.\allowbreak{}} & 12밴드 & \textbf{공백}으로 구분한 정수 nm. \texttt{produce\_\allowbreak{}grid.\allowbreak{}py}의 쉼표 형식과 다름 \\
\texttt{-\allowbreak{}-\allowbreak{}rows LO:\allowbreak{}HI} & 전체 & 0부터 세는 CSV 자료 행 슬라이스. LO 포함, HI 제외. 예: \texttt{0:\allowbreak{}10}은 첫 10행. 두 정수를 모두 입력 \\
\texttt{-\allowbreak{}-\allowbreak{}resume} & 끄기 & 기존 \texttt{(row\_\allowbreak{}id,\allowbreak{}wavelength\_\allowbreak{}nm)}를 건너뜀. 생략 시 기존 결과를 덮어씀 \\
\texttt{-\allowbreak{}-\allowbreak{}progress N} & 1 & N개 처리마다 진행 출력. 1 이상의 정수 \\
\end{longtable}
}

두 서브명령에는 \texttt{-\allowbreak{}-\allowbreak{}data} 옵션이 없다. 스크립트 옆 \texttt{data/\allowbreak{}}를 쓴다. \texttt{grid}의 선택적 slope 열은 해당 행에서 CLI 기본값보다 우선한다. slope 추가 별칭은 CDOM \texttt{ccrr\_\allowbreak{}adom\_\allowbreak{}slope/\allowbreak{}s\_\allowbreak{}cdom/\allowbreak{}scdom}, 쇄설물 \texttt{s\_\allowbreak{}detritus/\allowbreak{}sdetritus/\allowbreak{}s\_\allowbreak{}det}도 지원한다.

\texttt{grid}는 한 조건에서 예외가 발생하면 \texttt{FAIL row=\allowbreak{}.\allowbreak{}.\allowbreak{}.\allowbreak{} wl=\allowbreak{}.\allowbreak{}.\allowbreak{}.\allowbreak{}}을 출력하고 다음 조건으로 넘어간다. 마지막에 “완료”가 나와도 모든 조건이 성공했다는 뜻이 아니다. 요청한 행 수 \ensuremath{\times} 밴드 수와 실제 고유 결과 키 수를 비교하고 로그의 \texttt{FAIL}을 확인한다. 재개 시 헤더·설정·값 검증은 배치 생산기보다 약하므로 다른 버전의 CSV에 이어 쓰지 않는다.

\subsection{종 이름 목록}
\label{sec:auto:06_python_batch:5.1}
파서가 받는 호환 별칭은 \texttt{pico}, \texttt{nano}, \texttt{micro}이며, 정식 카탈로그 이름은 다음과 같다.


\begin{ManualCode}
eap_diatoms_pennate        eap_chlorophytes        eap_diatoms_centric
eap_cryptophytes           eap_cyano_blue         eap_cyano_red
eap_dinoflagellates        eap_eustigmatophytes   eap_hapto_pavlovaceae
eap_pelagophytes           eap_prasinophytes      eap_prochlorococcus
eap_hapto_prymnesiaceae     eap_raphidophytes      eap_rhodophytes
eap_synechococcus          eap_microcystis
\end{ManualCode}

이 목록이 생산 RT에서 종별 산란을 사용할 수 있다는 뜻은 아니다. Chl > 0에서는 명시적 선택 자체가 거절된다. Chl=0이라도 기본 \texttt{micro}/\texttt{eap\_\allowbreak{}diatoms\_\allowbreak{}centric} 외 선택은 \texttt{OCRT\_\allowbreak{}ADVANCED=\allowbreak{}1}을 요구하고 해당 EAP 파일을 읽는다. 일반 자료 생성에서는 종 옵션을 생략한다.

\section{\texttt{ocrt\_\allowbreak{}fullgrid.\allowbreak{}py}: 각도 LUT 전체 옵션}
\label{sec:auto:06_python_batch:6}
한 파장·태양각·대기·해양 상태에 대해 VZA와 RAA의 데카르트 곱을 계산한다. \texttt{produce\_\allowbreak{}grid.\allowbreak{}py}처럼 서로 다른 물리 조건을 입력 행마다 바꾸는 도구가 아니다.


{\TableFont
\begin{longtable}{@{}>{\raggedright\arraybackslash}p{\dimexpr 0.3000\linewidth-5.33pt\relax}>{\raggedright\arraybackslash}p{\dimexpr 0.2300\linewidth-5.33pt\relax}>{\raggedright\arraybackslash}p{\dimexpr 0.4700\linewidth-5.33pt\relax}@{}}
\toprule
\textbf{옵션} & \textbf{기본값} & \textbf{설명} \\
\midrule
\endfirsthead
\toprule
\textbf{옵션} & \textbf{기본값} & \textbf{설명} \\
\midrule
\endhead
\bottomrule
\endfoot
\texttt{-\allowbreak{}h}, \texttt{-\allowbreak{}-\allowbreak{}help} & --- & 도움말 \\
\texttt{-\allowbreak{}-\allowbreak{}out PATH} & 필수 & 각도 LUT CSV. 기존 파일 덮어씀, 자동 재개 없음 \\
\texttt{-\allowbreak{}-\allowbreak{}wl NM}, \texttt{-\allowbreak{}-\allowbreak{}sza DEG}, \texttt{-\allowbreak{}-\allowbreak{}wind MPS} & 각각 필수 & 파장·태양각·풍속 \\
\texttt{-\allowbreak{}-\allowbreak{}chl}, \texttt{-\allowbreak{}-\allowbreak{}tsm}, \texttt{-\allowbreak{}-\allowbreak{}acdom440}, \texttt{-\allowbreak{}-\allowbreak{}aod865} & 각각 필수 & 값이 0이어도 명시 \\
\texttt{-\allowbreak{}-\allowbreak{}aer NAME} & \texttt{C50} & 에어로졸 모델명/경로, 확장자 선택 가능 \\
\texttt{-\allowbreak{}-\allowbreak{}vza-\allowbreak{}values LIST} & \texttt{0,\allowbreak{}30} & 쉼표로 구분한 실수 각도 목록 \\
\texttt{-\allowbreak{}-\allowbreak{}raa-\allowbreak{}values LIST} & \texttt{0,\allowbreak{}90,\allowbreak{}180,\allowbreak{}270} & 쉼표로 구분한 OCRT RAA 목록 \\
\texttt{-\allowbreak{}-\allowbreak{}phyto-\allowbreak{}group NAME} & 생략 & \ref{sec:auto:06_python_batch:5}절의 종 제한 \\
\texttt{-\allowbreak{}-\allowbreak{}tsm-\allowbreak{}species NAME} & \texttt{red\_\allowbreak{}clay} & 광물 네 종류 중 선택 \\
\texttt{-\allowbreak{}-\allowbreak{}cdom-\allowbreak{}slope S} & 0.014 & CDOM 흡수 기울기, nm\textsuperscript{-}\textsuperscript{1} \\
\texttt{-\allowbreak{}-\allowbreak{}n-\allowbreak{}mu-\allowbreak{}water N} & 48 & 결합 계산 Gauss 절점 수 \\
\texttt{-\allowbreak{}-\allowbreak{}m-\allowbreak{}max N} & 16 & 대기 Fourier 상한 \\
\texttt{-\allowbreak{}-\allowbreak{}n-\allowbreak{}layers-\allowbreak{}atm N} & 400 & 대기 층 수 \\
\texttt{-\allowbreak{}-\allowbreak{}water-\allowbreak{}m-\allowbreak{}max N} & 30 & 수중 Fourier 상한 \\
\texttt{-\allowbreak{}-\allowbreak{}water-\allowbreak{}n-\allowbreak{}layers N} & 300 & 수중 층 수 설정. 실제 수중 정책에서 필요한 층 수가 더 클 수 있음 \\
\texttt{-\allowbreak{}-\allowbreak{}max-\allowbreak{}it-\allowbreak{}water N} & 500 & 수중 산란차수 상한 \\
\texttt{-\allowbreak{}-\allowbreak{}pressure HPA} & 1013.25 & 압력, hPa \\
\texttt{-\allowbreak{}-\allowbreak{}no-\allowbreak{}gas} & 미지정 & 기체흡수 끄기 \\
\end{longtable}
}

\texttt{-\allowbreak{}-\allowbreak{}data}, \texttt{-\allowbreak{}-\allowbreak{}which}, \texttt{-\allowbreak{}-\allowbreak{}detritus-\allowbreak{}a440}, \texttt{-\allowbreak{}-\allowbreak{}detritus-\allowbreak{}slope}는 이 CLI에 없다. 해수 굴절률은 1.34이다. Mie와 파장별 에어로졸 상태를 한 번 준비해 각도 격자에서 공유한다.

\section{실행 예시}
\label{sec:auto:06_python_batch:7}
\textbf{아래 R1 예시는 \ref{sec:auto:06_python_batch:2}절의 성분 자료 준비 검사가 성공한 환경에서 실행한다. 현재 로컬에서는 EAP 누락으로 R1의 수치 실행을 검증하지 못했다.} 명령 구문과 옵션은 실제 파서·소스에 대조했다. 대규모 실행 전에 실제 자료·수치 조건으로 소규모 대조를 수행한다.

\subsection{Rayleigh 기준 반사도 한 점}
\label{sec:auto:06_python_batch:7.1}
에어로졸·기체흡수 없는 거친 Fresnel 해면 위 Rayleigh 기준이다. 완전 흑색 경계의 대기와는 해면 조건이 다르다.


\begin{ManualCode}
python ocrt_solve.py single --wl 555 --sza 30 --vza 20 --raa 90 --wind 5 --aer C50 --aod865 0 --which R2 --no-gas --pressure 1013.25 --n-layers 400
\end{ManualCode}

R2는 해수 성분 모델을 만들지 않으므로 EAP 파일이 필요하지 않지만 \texttt{C50.\allowbreak{}mie}는 읽는다. 파서에 지정하지 않은 R1/R3·Rrs 출력은 없다. IPSS를 포함한 운영용 해양 대기보정 Rayleigh LUT는 C 예제를 사용한다.

매뉴얼 작성 시 이 명령을 Windows CPU/NumPy 2.3.5에서 실제 실행해 정상 종료와 \texttt{rho\_\allowbreak{}R\_\allowbreak{}I =\allowbreak{} 3.\allowbreak{}8982239255e-\allowbreak{}02}, \texttt{aod\_\allowbreak{}band =\allowbreak{} 0}을 확인했다. 이 값은 해당 설치의 동작 확인 기록이며, C와의 정밀도 검증 기준을 새로 확정한 값은 아니다.

\subsection{여러 압력·파장의 소규모 Rayleigh LUT}
\label{sec:auto:06_python_batch:7.2}
다음 코드는 CSV를 생성하는 예다. \texttt{python} 대화형 입력 또는 별도 파일로 실행한다.


\begin{ManualCode}
import csv
from itertools import product
from pathlib import Path

Path("results").mkdir(exist_ok=True)
with open("results/rayleigh_grid.csv", "w", newline="", encoding="utf-8") as f:
    w = csv.writer(f)
    w.writerow(["case_id", "sza", "vza", "raa", "wind", "aerosol", "aod865", "chl", "tsm", "acdom440"])
    for i, (s, v, r, wind) in enumerate(product([20, 40, 60], [0.001, 20, 40], [0, 60, 120, 180, 240, 300], [2, 5, 10])):
        w.writerow([f"ray_{i:04d}", s, v, r, wind, "C50", 0, 0, 0, 0])
\end{ManualCode}

각 압력에서 162조건 \ensuremath{\times} 3파장 = 486행을 기대한다.


\begin{ManualCode}
python ocrt_solve.py grid --grid results/rayleigh_grid.csv --out results/rayleigh_p1013.csv --bands 443 555 865 --which R2 --no-gas --pressure 1013.25 --n-layers 400
python ocrt_solve.py grid --grid results/rayleigh_grid.csv --out results/rayleigh_p900.csv --bands 443 555 865 --which R2 --no-gas --pressure 900 --n-layers 400
\end{ManualCode}

이 LUT의 제품은 \texttt{rho\_\allowbreak{}R\_\allowbreak{}I/\allowbreak{}Q/\allowbreak{}U}이며, 압력은 출력 CSV 열에 자동 저장되지 않으므로 실행 명령과 함께 보관한다. 중단 후에는 같은 명령에 \texttt{-\allowbreak{}-\allowbreak{}resume}을 붙인다. 태양·관측 기하, 풍속, 기체 적용 여부, 압력별 보간법을 LUT 메타데이터에 남긴다. 이 예는 Python 평면 대기 경로의 검토용 LUT이다.

\subsection{대양·연안·탁수·방위 변화 시뮬레이션}
\label{sec:auto:06_python_batch:7.3}
문서에 포함된 \href{https://github.com/brtnt/OCRT/blob/main/docs/manual/examples/python_simulation_grid.csv}{python\_simulation\_grid.csv}는 6개 조건이다. 해수 상태·AOD·풍속·방위·흡수 기울기를 명시한 사용 예이며, 검증된 자연환경 분포를 의미하지 않는다.


\begin{ManualCode}
New-Item -ItemType Directory -Force results | Out-Null
python ocrt_solve.py grid --grid ../../docs/manual/examples/python_simulation_grid.csv --out results/simulation_single_443_555.csv --bands 443 555 --n-mu-water 24 --n-layers 400 --m-max 16 --max-it-water 500
python produce_grid.py --grid ../../docs/manual/examples/python_simulation_grid.csv --out results/simulation_batch_443_555.csv --bands 443,555 --chunk 2 --n-mu-water 24 --nt-atm 400 --m-max 16 --max-it-water 500
\end{ManualCode}

각 파일에서 12개 고유 조건·파장 키를 기대한다. 단일과 배치의 공통 \texttt{rho\_\allowbreak{}TOA}, R2/R3, Rrs를 맞춰 비교하고 차이를 기록한다. 위의 절점 수를 24로 맞춰도 대기 커널과 R2 설정은 같아지지 않는다. 일치가 필요한 생산 목적에서는 그 차이부터 해소·검증해야 한다. 낮은 산란차수나 층 수로 실행한 시험 결과를 최종 수렴 결과로 사용하지 않는다.

탁수 행은 특히 산란차수 상한만으로 수렴 여부를 판단하지 않는다. \texttt{ocrt\_\allowbreak{}solve.\allowbreak{}py grid}의 \texttt{conv/\allowbreak{}orders}를 확인하고 \texttt{-\allowbreak{}-\allowbreak{}n-\allowbreak{}mu-\allowbreak{}water}, \texttt{-\allowbreak{}-\allowbreak{}n-\allowbreak{}layers}, \texttt{-\allowbreak{}-\allowbreak{}m-\allowbreak{}max}, \texttt{-\allowbreak{}-\allowbreak{}max-\allowbreak{}it-\allowbreak{}water}를 늘린 별도 파일과 필요한 제품의 변화를 비교한다. \texttt{produce\_\allowbreak{}grid.\allowbreak{}py} 출력에는 수중 수렴 플래그가 없으므로 별도의 대표 조건 검사가 필요하다.

\subsection{기체·CDOM·광물 민감도}
\label{sec:auto:06_python_batch:7.4}
기체 효과를 비교하려면 동일 CSV를 새 출력 파일로 실행한다.


\begin{ManualCode}
python produce_grid.py --grid ../../docs/manual/examples/python_simulation_grid.csv --out results/simulation_nogas.csv --bands 443,555 --chunk 2 --no-gas
\end{ManualCode}

이 비교에서는 앞선 실행의 수치 옵션도 동일하게 지정한다. 기체 제거가 물 IOP를 제거하지는 않는다. CDOM 기울기만 비교할 때는 \texttt{single -\allowbreak{}-\allowbreak{}adom-\allowbreak{}slope 0.\allowbreak{}014}와 \texttt{single -\allowbreak{}-\allowbreak{}adom-\allowbreak{}slope 0.\allowbreak{}018}을 쓰거나 CSV에 서로 다른 \texttt{ocrt\_\allowbreak{}adom\_\allowbreak{}slope}를 명시한다. \texttt{grid}에서 CSV slope 열이 채워져 있으면 CLI 값을 바꿔도 해당 행의 값이 우선한다.

광물 종류는 \texttt{ocrt\_\allowbreak{}solve.\allowbreak{}py -\allowbreak{}-\allowbreak{}tsm-\allowbreak{}species} 또는 \texttt{produce\_\allowbreak{}grid.\allowbreak{}py -\allowbreak{}-\allowbreak{}mineral}로 바꾸며 각 실행을 별도 파일에 기록한다. 같은 TSM 질량농도라도 종류가 다르면 IOP가 달라진다.

\subsection{각도별 수출광·편광 LUT}
\label{sec:auto:06_python_batch:7.5}

\begin{ManualCode}
python ocrt_fullgrid.py --out results/angular_555.csv --wl 555 --sza 30 --wind 5 --chl 1 --tsm 3 --acdom440 0.04 --aod865 0.1 --aer C50 --vza-values 0.001,20,40,55 --raa-values 0,45,90,135,180,225,270,315 --n-mu-water 48 --m-max 16 --n-layers-atm 400 --water-m-max 30 --water-n-layers 300 --max-it-water 500
\end{ManualCode}

32행의 \texttt{Rrs\_\allowbreak{}I/\allowbreak{}Q/\allowbreak{}U}, \texttt{rrs\_\allowbreak{}I/\allowbreak{}Q/\allowbreak{}U}, \texttt{TOA\_\allowbreak{}rho\_\allowbreak{}I/\allowbreak{}Q/\allowbreak{}U}를 얻는다. \texttt{water\_\allowbreak{}converged}와 \texttt{water\_\allowbreak{}orders}를 확인한다. 이 CSV는 \texttt{produce\_\allowbreak{}grid.\allowbreak{}py}와 열 구조가 다르다. 물리 상태마다 다른 파일을 만들어 파장·SZA·AOD·Chl·TSM·CDOM·풍속 차원의 LUT로 묶을 수 있다. 단일 각도 출력과 공유 각도장을 이용한 LUT 출력의 절점·수렴 차이도 확인한다.

\subsection{재현 가능한 무작위 입력 설계와 소규모 추출}
\label{sec:auto:06_python_batch:7.6}

\begin{ManualCode}
python make_grid.py --n 1000 --seed 20260718 --out results/design_1000.csv --iop-out results/design_water_1000.csv
python subsample_grid.py --in results/design_1000.csv --n 8 --seed 20260720 --out results/pilot_8.csv
\end{ManualCode}

\texttt{make\_\allowbreak{}grid.\allowbreak{}py}는 자연로그 공간의 상관 정규분포를 범위 절단해서 Chl·TSM·CDOM을 뽑는다. TSM \texttt{0.\allowbreak{}1--20 g/\allowbreak{}m\textsuperscript{3}}, Chl \texttt{0.\allowbreak{}1--5 mg/\allowbreak{}m\textsuperscript{3}}, CDOM \texttt{0.\allowbreak{}01--0.\allowbreak{}1 m\textsuperscript{-}\textsuperscript{1}}; 목표 로그상관은 TSM--Chl \texttt{0.\allowbreak{}85}, Chl--CDOM \texttt{0.\allowbreak{}75}, TSM--CDOM \texttt{0.\allowbreak{}77}이다. AOD865는 \texttt{0.\allowbreak{}05--0.\allowbreak{}3}, 풍속은 \texttt{1--10 m/\allowbreak{}s}, SZA는 \texttt{0--65\textdegree{}}, VZA는 \texttt{0--55\textdegree{}}, RAA는 \texttt{0--180\textdegree{}}; 에어로졸은 \texttt{C50/\allowbreak{}T50/\allowbreak{}M80C/\allowbreak{}M95C/\allowbreak{}M98C/\allowbreak{}O99} 중 무작위 선택이다. 직접 글린트 반사도 절댓값 \texttt{<5e-\allowbreak{}4}를 만족할 때까지 기하를 다시 뽑는다. 따라서 채택된 기하와 풍속이 완전히 독립인 표본은 아니다.

범위 절단 뒤의 상관은 목표와 달라질 수 있다. 원래 논문용 과거 격자와 같은 행들을 복원하는 도구도 아니다. 생성 CSV와 시드를 함께 보관한다.


{\TableFont
\begin{longtable}{@{}>{\raggedright\arraybackslash}p{\dimexpr 0.3400\linewidth-4.00pt\relax}>{\raggedright\arraybackslash}p{\dimexpr 0.6600\linewidth-4.00pt\relax}@{}}
\toprule
\textbf{도구} & \textbf{전체 옵션} \\
\midrule
\endfirsthead
\toprule
\textbf{도구} & \textbf{전체 옵션} \\
\midrule
\endhead
\bottomrule
\endfoot
\texttt{make\_\allowbreak{}grid.\allowbreak{}py} & \texttt{-\allowbreak{}h/\allowbreak{}-\allowbreak{}-\allowbreak{}help}; \texttt{-\allowbreak{}-\allowbreak{}n} 기본 1000; \texttt{-\allowbreak{}-\allowbreak{}seed} 기본 20260718; \texttt{-\allowbreak{}-\allowbreak{}out} 기본 \texttt{full\_\allowbreak{}grid\_\allowbreak{}design\_\allowbreak{}v1.\allowbreak{}csv}; \texttt{-\allowbreak{}-\allowbreak{}iop-\allowbreak{}out} 기본 \texttt{iop\_\allowbreak{}grid\_\allowbreak{}design\_\allowbreak{}v1.\allowbreak{}csv} \\
\texttt{subsample\_\allowbreak{}grid.\allowbreak{}py} & \texttt{-\allowbreak{}h/\allowbreak{}-\allowbreak{}-\allowbreak{}help}; \texttt{-\allowbreak{}-\allowbreak{}in} 기본 \texttt{full\_\allowbreak{}grid\_\allowbreak{}design\_\allowbreak{}v1.\allowbreak{}csv}; \texttt{-\allowbreak{}-\allowbreak{}n} 기본 100; \texttt{-\allowbreak{}-\allowbreak{}seed} 기본 20260720; \texttt{-\allowbreak{}-\allowbreak{}out} 생략 시 \texttt{draft\_\allowbreak{}grid\_\allowbreak{}<실제 행 수>.\allowbreak{}csv} \\
\end{longtable}
}

부분집합 추출은 원본 ID와 원본 행 순서를 유지한다. 요청 수가 원본보다 크면 전체 행을 선택한다. 현재 요약 출력에서 첫 두 행을 참조하므로 \texttt{-\allowbreak{}-\allowbreak{}n 0}/\texttt{-\allowbreak{}-\allowbreak{}n 1} 대신 2행 이상을 사용한다.

\Needspace{7\baselineskip}
\section{출력 CSV 해석}
\label{sec:auto:06_python_batch:8}
\subsection{\texttt{produce\_\allowbreak{}grid.\allowbreak{}py}의 열}
\label{sec:auto:06_python_batch:8.1}
반사도 \texttt{rho\_\allowbreak{}*}는 무차원이다. 아래 차감 제품은 해당 실행에서 계산한 세 반사도의 대수적 차이다. \textbf{차감 항등식이 성립한다는 사실만으로 물리 성분 분리가 정확하거나 C와 일치한다고 판단할 수 없다.} 서로 다른 수치 커널·절점·산란차수의 영향도 차이에 남을 수 있다.


{\TableFont
\begin{longtable}{@{}>{\raggedright\arraybackslash}p{\dimexpr 0.3400\linewidth-4.00pt\relax}>{\raggedright\arraybackslash}p{\dimexpr 0.6600\linewidth-4.00pt\relax}@{}}
\toprule
\textbf{열} & \textbf{의미} \\
\midrule
\endfirsthead
\toprule
\textbf{열} & \textbf{의미} \\
\midrule
\endhead
\bottomrule
\endfoot
\texttt{case\_\allowbreak{}id}, \texttt{band\_\allowbreak{}nm} & 조건 ID와 계산 파장 nm \\
\texttt{sza,\allowbreak{}vza,\allowbreak{}raa,\allowbreak{}wind,\allowbreak{}aerosol,\allowbreak{}aod865,\allowbreak{}chl,\allowbreak{}tsm,\allowbreak{}acdom440} & 입력 조건 복사. 각도·풍속은 소수 4자리, 농도/AOD는 소수 6자리로 저장 \\
\texttt{rho\_\allowbreak{}TOA} & R1: 대기--해양 결합 I 반사도. 직접 선글린트를 분리한 제품 \\
\texttt{rho\_\allowbreak{}R} & R2: 에어로졸 없는 Rayleigh + 거친 Fresnel 해면 기준 I 반사도 \\
\texttt{rho\_\allowbreak{}RpA} & R3: Rayleigh + 에어로졸 + 흑색 해양 위 Fresnel 해면 I 반사도 \\
\texttt{rho\_\allowbreak{}RC} & \texttt{rho\_\allowbreak{}TOA - rho\_\allowbreak{}R} \\
\texttt{rho\_\allowbreak{}A\_\allowbreak{}RA} & \texttt{rho\_\allowbreak{}RpA - rho\_\allowbreak{}R} \\
\texttt{t\_\allowbreak{}rho\_\allowbreak{}w} & \texttt{rho\_\allowbreak{}TOA - rho\_\allowbreak{}RpA} \\
\texttt{T\_\allowbreak{}dir\_\allowbreak{}dn} & 태양 방향 직달 대기투과율 \texttt{exp(-\ensuremath{\tau}/\allowbreak{}\ensuremath{\mu}sun)} \\
\texttt{T\_\allowbreak{}diff\_\allowbreak{}dn\_\allowbreak{}hemi} & 하향 확산광의 반구 적분 기여 \\
\texttt{T\_\allowbreak{}total\_\allowbreak{}dn\_\allowbreak{}hemi} & \texttt{Ed(0+) /\allowbreak{} (Fsun·\ensuremath{\mu}sun)} \\
\texttt{T\_\allowbreak{}dir\_\allowbreak{}up\_\allowbreak{}view} & 관측 방향 직달 대기투과율 \texttt{exp(-\ensuremath{\tau}/\allowbreak{}\ensuremath{\mu}view)} \\
\texttt{T\_\allowbreak{}total\_\allowbreak{}up\_\allowbreak{}view} & 실제 bottom-source 패스의 TOA 수광 radiance를 \texttt{Lu(0+)}로 나눈 유효 방향 투과율 \\
\texttt{T\_\allowbreak{}diff\_\allowbreak{}up\_\allowbreak{}view} & \texttt{T\_\allowbreak{}total\_\allowbreak{}up\_\allowbreak{}view - T\_\allowbreak{}dir\_\allowbreak{}up\_\allowbreak{}view}. 각도 재분배 때문에 음수 가능; 일률적으로 0으로 잘라내지 않음 \\
\texttt{TOA\_\allowbreak{}water\_\allowbreak{}signal\_\allowbreak{}I} & R1 내부 수중 기원 광원을 대기로 전달한 TOA radiance. \texttt{t\_\allowbreak{}rho\_\allowbreak{}w}와 단위가 다름 \\
\texttt{T\_\allowbreak{}up\_\allowbreak{}rt\_\allowbreak{}valid} & \texttt{abs(Lu(0+)) > 0}이면 1, 아니면 0. 수치 수렴이나 전체 품질 판정 플래그가 아님 \\
\texttt{rho\_\allowbreak{}wn} & \texttt{t\_\allowbreak{}rho\_\allowbreak{}w /\allowbreak{} (T\_\allowbreak{}total\_\allowbreak{}dn\_\allowbreak{}hemi·T\_\allowbreak{}total\_\allowbreak{}up\_\allowbreak{}view)}. 분모 0/NaN이면 NaN. 정규화 정의와 R1/R3 차이를 고려해 사용 \\
\texttt{Rrs}, \texttt{Rrs\_\allowbreak{}I} & 같은 값. \texttt{Lu\_\allowbreak{}I(0+)/\allowbreak{}Ed(0+)}, sr\textsuperscript{-}\textsuperscript{1} \\
\texttt{Rrs\_\allowbreak{}Q}, \texttt{Rrs\_\allowbreak{}U} & \texttt{Lu\_\allowbreak{}Q/\allowbreak{}U(0+)/\allowbreak{}Ed(0+)}, sr\textsuperscript{-}\textsuperscript{1}. I/Q/U 모두 동일한 I 하향 irradiance로 나눔 \\
\texttt{rrs\_\allowbreak{}I}, \texttt{rrs\_\allowbreak{}Q}, \texttt{rrs\_\allowbreak{}U} & 수면 바로 아래 \texttt{Lu\_\allowbreak{}I/\allowbreak{}Q/\allowbreak{}U(0-)/\allowbreak{}Ed(0-)}, sr\textsuperscript{-}\textsuperscript{1} \\
\texttt{a\_\allowbreak{}total,\allowbreak{}b\_\allowbreak{}total,\allowbreak{}bb\_\allowbreak{}total} & 해수 총 흡수·산란·후방산란계수, m\textsuperscript{-}\textsuperscript{1} \\
\texttt{a\_\allowbreak{}w,\allowbreak{}b\_\allowbreak{}w,\allowbreak{}bb\_\allowbreak{}w} & 순수해수의 같은 세 계수, m\textsuperscript{-}\textsuperscript{1} \\
\texttt{a\_\allowbreak{}chl} & 식물플랑크톤 흡수 기여만, m\textsuperscript{-}\textsuperscript{1} \\
\texttt{a\_\allowbreak{}phyto\_\allowbreak{}detritus,\allowbreak{}b\_\allowbreak{}phyto\_\allowbreak{}detritus,\allowbreak{}bb\_\allowbreak{}phyto\_\allowbreak{}detritus} & 식물플랑크톤 + Chl 연동 쇄설물 합, m\textsuperscript{-}\textsuperscript{1}. 현재 식물플랑크톤 산란은 0 \\
\texttt{a\_\allowbreak{}dom} & CDOM 흡수, m\textsuperscript{-}\textsuperscript{1} \\
\texttt{a\_\allowbreak{}min,\allowbreak{}b\_\allowbreak{}min,\allowbreak{}bb\_\allowbreak{}min} & 광물의 흡수·산란·후방산란, m\textsuperscript{-}\textsuperscript{1} \\
\texttt{Kd0minus} & 수면 바로 아래 하향 irradiance의 감쇠계수, m\textsuperscript{-}\textsuperscript{1} \\
\end{longtable}
}

\texttt{T\_\allowbreak{}up\_\allowbreak{}rt\_\allowbreak{}valid=\allowbreak{}0}인 행의 상향 유효투과율과 \texttt{rho\_\allowbreak{}wn}은 NaN일 수 있다. NaN을 0으로 대체해서 학습·보정 자료에 넣지 않는다. 위 CSV에는 TOA Q/U, slope 선택값, 광물명, 압력, 버전, GPU/CPU 모드, 수치 설정, 수중 수렴 상태가 모두 저장되는 것은 아니다. 특히 optional slope 열은 입력에는 있어도 출력에 그대로 복사되지 않는다. 원본 CSV와 명령 로그를 반드시 함께 보관한다.

Rrs가 필요하면 \texttt{Rrs\_\allowbreak{}I}를 직접 읽는다. \texttt{rho\_\allowbreak{}wn/\allowbreak{}\ensuremath{\pi}}와 Rrs의 관계를 검증하지 않고 서로 대체하지 않는다. \texttt{Rrs\_\allowbreak{}Q/\allowbreak{}U}가 음수라는 이유로 불량값으로 제거해서도 안 된다.

\subsection{\texttt{ocrt\_\allowbreak{}solve.\allowbreak{}py}의 출력}
\label{sec:auto:06_python_batch:8.2}
\texttt{single}은 CSV가 아니라 표준출력 텍스트를 낸다. \texttt{conv/\allowbreak{}orders}와 일부 반사도·Rrs를 보여 준다. \texttt{-\allowbreak{}-\allowbreak{}which R2}처럼 R1을 실행하지 않은 경우 \texttt{conv/\allowbreak{}orders}는 \texttt{None}이다. R1에서 Q/U Rrs를 내부적으로 계산하더라도 현재 이 CLI 텍스트는 그 모든 키를 출력하지 않는다.

\texttt{grid}는 다음 열들을 저장한다. 요청하지 않은 반사도 또는 조합에 필요한 피연산자가 없으면 해당 필드는 빈칸이다.


\begin{ManualCode}
row_id,wavelength_nm,tsm,chl,acdom440,aer_model,aod865,wind,sza,vza,raa,aod_band,
rho_TOA_I,rho_R_I,rho_RA_black_I,rho_RC_I,rho_AplusRA_I,t_rho_w_I,Rrs,rrs0minus,
rho_TOA_Q,rho_TOA_U,rho_R_Q,rho_R_U,rho_RA_black_Q,rho_RA_black_U,orders,conv,
adom_slope,detritus_a440,detritus_slope
\end{ManualCode}

\texttt{rho\_\allowbreak{}RA\_\allowbreak{}black\_\allowbreak{}*}가 R3이고, \texttt{rho\_\allowbreak{}AplusRA\_\allowbreak{}I}는 R3-R2이다. \texttt{aod\_\allowbreak{}band}는 Mie 분광 소광비로 865 nm AOD를 계산 파장으로 변환한 값이다. \texttt{Rrs/\allowbreak{}rrs0minus}는 I 성분의 수면 위/아래 비율이다. 수중 \texttt{conv=\allowbreak{}1}을 확인하며, 이는 대기·전체 모델의 모든 정확성 검사를 대체하지 않는다. \texttt{orders}는 사용한 수중 최대 산란차수다.

\subsection{\texttt{ocrt\_\allowbreak{}fullgrid.\allowbreak{}py}의 출력}
\label{sec:auto:06_python_batch:8.3}

{\TableFont
\begin{longtable}{@{}>{\raggedright\arraybackslash}p{\dimexpr 0.3400\linewidth-4.00pt\relax}>{\raggedright\arraybackslash}p{\dimexpr 0.6600\linewidth-4.00pt\relax}@{}}
\toprule
\textbf{열 묶음} & \textbf{의미} \\
\midrule
\endfirsthead
\toprule
\textbf{열 묶음} & \textbf{의미} \\
\midrule
\endhead
\bottomrule
\endfoot
\texttt{case\_\allowbreak{}kind} & 현재 \texttt{ocean\_\allowbreak{}rrs\_\allowbreak{}grid} \\
\texttt{sza\_\allowbreak{}deg,\allowbreak{}wavelength\_\allowbreak{}nm,\allowbreak{}wind\_\allowbreak{}speed\_\allowbreak{}ms} & 물리 조건 \\
\texttt{mu\_\allowbreak{}index,\allowbreak{}mu\_\allowbreak{}view,\allowbreak{}vza\_\allowbreak{}deg,\allowbreak{}raa\_\allowbreak{}deg} & VZA 목록의 0기반 인덱스, \texttt{cos(VZA)}, 입력 각도 \\
\texttt{Ed0plus\_\allowbreak{}air,\allowbreak{}Lu0plus\_\allowbreak{}I,\allowbreak{}Lu0plus\_\allowbreak{}Q,\allowbreak{}Lu0plus\_\allowbreak{}U} & 수면 위 하향 irradiance와 수출광 Stokes radiance \\
\texttt{Rrs\_\allowbreak{}I,\allowbreak{}Rrs\_\allowbreak{}Q,\allowbreak{}Rrs\_\allowbreak{}U} & 위 radiance를 \texttt{Ed0plus\_\allowbreak{}air}로 나눈 값, sr\textsuperscript{-}\textsuperscript{1} \\
\texttt{Ed0minus\_\allowbreak{}water,\allowbreak{}Lu0minus\_\allowbreak{}I,\allowbreak{}Lu0minus\_\allowbreak{}Q,\allowbreak{}Lu0minus\_\allowbreak{}U} & 수면 아래 대응 값 \\
\texttt{rrs\_\allowbreak{}I,\allowbreak{}rrs\_\allowbreak{}Q,\allowbreak{}rrs\_\allowbreak{}U} & 수면 아래 radiance/irradiance 비, sr\textsuperscript{-}\textsuperscript{1} \\
\texttt{TOA\_\allowbreak{}rho\_\allowbreak{}I,\allowbreak{}TOA\_\allowbreak{}rho\_\allowbreak{}Q,\allowbreak{}TOA\_\allowbreak{}rho\_\allowbreak{}U} & 직접 선글린트를 분리한 TOA Stokes 반사도 \\
\texttt{a\_\allowbreak{}total,\allowbreak{}b\_\allowbreak{}total,\allowbreak{}bb\_\allowbreak{}total,\allowbreak{}omega\_\allowbreak{}water} & 해수 IOP 세 계수와 \texttt{b/\allowbreak{}(a+b)} \\
\texttt{water\_\allowbreak{}orders,\allowbreak{}water\_\allowbreak{}converged} & 수중 최대 사용 산란차수·수렴 플래그 \\
\texttt{aod\_\allowbreak{}ref,\allowbreak{}aod\_\allowbreak{}ref\_\allowbreak{}nm,\allowbreak{}aod\_\allowbreak{}target} & 기준 AOD·기준 파장 nm·계산 파장의 AOD \\
\texttt{tau\_\allowbreak{}a\_\allowbreak{}eff,\allowbreak{}ssa\_\allowbreak{}a\_\allowbreak{}eff} & 에어로졸 처리 후 유효 광학두께·단일산란알베도 \\
\end{longtable}
}

irradiance/radiance는 내부 \texttt{F\_\allowbreak{}sun\_\allowbreak{}TOA=\allowbreak{}1} 정규화에 따른 값이며, 실제 관측의 W 단위 분광복사휘도에 바로 대응시키지 않는다. 물리 절대량이 필요하면 입사 태양 스펙트럼의 정규화를 일관되게 적용한다. CSV에는 Chl/TSM/CDOM 입력 자체가 없으므로 상태 파일과 함께 묶는다.

\section{GPU 대조와 검증 도구}
\label{sec:auto:06_python_batch:9}

\begin{ManualCode}
python gpu_gate.py --grid results/pilot_8.csv --data data --bands 443,555 --chunk 4 --n-mu-water 24 --nt-atm 400 --m-max 16 --max-it-water 500 --tol 1e-7
\end{ManualCode}


{\TableFont
\begin{longtable}{@{}>{\raggedright\arraybackslash}p{\dimexpr 0.3000\linewidth-5.33pt\relax}>{\raggedright\arraybackslash}p{\dimexpr 0.2300\linewidth-5.33pt\relax}>{\raggedright\arraybackslash}p{\dimexpr 0.4700\linewidth-5.33pt\relax}@{}}
\toprule
\textbf{옵션} & \textbf{기본값} & \textbf{설명} \\
\midrule
\endfirsthead
\toprule
\textbf{옵션} & \textbf{기본값} & \textbf{설명} \\
\midrule
\endhead
\bottomrule
\endfoot
\texttt{-\allowbreak{}h}, \texttt{-\allowbreak{}-\allowbreak{}help} & --- & 도움말 \\
\texttt{-\allowbreak{}-\allowbreak{}grid PATH} & 필수 & 소규모 조건 CSV \\
\texttt{-\allowbreak{}-\allowbreak{}data DIR} & 현재 폴더의 \texttt{data} & 자료 루트. \texttt{produce\_\allowbreak{}grid.\allowbreak{}py}와 기본 경로 기준이 다름 \\
\texttt{-\allowbreak{}-\allowbreak{}bands LIST} & \texttt{443,\allowbreak{}555,\allowbreak{}865} & 쉼표로 구분한 정수 파장 \\
\texttt{-\allowbreak{}-\allowbreak{}chunk N} & 8 & 배치 행 수 \\
\texttt{-\allowbreak{}-\allowbreak{}n-\allowbreak{}mu-\allowbreak{}water N} & 24 & 수중/배치 대기 적분 절점 수 \\
\texttt{-\allowbreak{}-\allowbreak{}nt-\allowbreak{}atm N} & 400 & R1/R3 대기 층 수 \\
\texttt{-\allowbreak{}-\allowbreak{}m-\allowbreak{}max N} & 16 & R1/R3 대기 Fourier 상한 \\
\texttt{-\allowbreak{}-\allowbreak{}max-\allowbreak{}it-\allowbreak{}water N} & 500 & 수중 산란차수 상한 \\
\texttt{-\allowbreak{}-\allowbreak{}tol VALUE} & \texttt{1e-\allowbreak{}7} & CPU 대비 상대차의 통과 기준. 계산기 자체 수렴오차 옵션이 아님 \\
\texttt{-\allowbreak{}-\allowbreak{}gpu-\allowbreak{}only} & 끄기 & CPU 비교를 생략하고 GPU 시간만 측정. 검증 통과로 해석하지 않음 \\
\end{longtable}
}

대조 대상은 \texttt{rho\_\allowbreak{}TOA/\allowbreak{}rho\_\allowbreak{}R/\allowbreak{}rho\_\allowbreak{}RpA/\allowbreak{}Rrs} 네 가지 I 제품이다. 기체흡수는 켜고, 기본 광물은 \texttt{red\_\allowbreak{}clay}, 기본 Chl 경로를 사용한다. 광물·기체 on/off를 바꾸는 옵션은 없다. Q/U, 모든 IOP·투과율, C 결과를 검사하지 않는다. 코드상 NaN·누락 키를 엄격히 실패 처리하는 검사가 충분하지 않으므로, \textbf{PASS 문구와 별개로 두 계산의 키 수·모든 대상 값의 유한성을 확인한다.} 기준값이 0이면 분모를 1로 놓아 절대차처럼 계산한다.

\texttt{verify.\allowbreak{}py}는 옵션 없는 과거 고정 기준값 검사다. 코드에 박힌 축소조건 C 기준 I 반사도 \texttt{2.\allowbreak{}2748060400e-\allowbreak{}01}와 상대차 \texttt{<1e-\allowbreak{}8}, 수중 \texttt{conv=\allowbreak{}1}을 검사한다. 현재 C 빌드·자료 SHA와 동기화된 기준임을 별도로 확인해야 하며, 이 한 케이스 PASS로 현재 C 전체 지원 범위를 보증하지 않는다. \texttt{diag\_\allowbreak{}gpu.\allowbreak{}py}는 옵션 없는 GPU 환경 진단 스크립트다.

고급 환경변수는 기본 생산 예제에서 설정하지 않는다. \texttt{OCRT\_\allowbreak{}PY\_\allowbreak{}NATIVE\_\allowbreak{}COUPLED\_\allowbreak{}LUT\_\allowbreak{}OFF}는 공유 각도 LUT 대신 과거 셀 반복 경로, \texttt{OCRT\_\allowbreak{}PY\_\allowbreak{}AIR\_\allowbreak{}SURFACE\_\allowbreak{}CACHE\_\allowbreak{}OFF}는 해면 커널 캐시 비활성화용이며 비어 있지 않은 값이면 작동한다. \texttt{OCRT\_\allowbreak{}PY\_\allowbreak{}SEQ\_\allowbreak{}BCS=\allowbreak{}1}은 배치 경계 적분의 순차 경로 대조용이다. \texttt{OCRT\_\allowbreak{}VALUE\_\allowbreak{}PHASE\_\allowbreak{}SPLINE}, \texttt{OCRT\_\allowbreak{}WATER\_\allowbreak{}VALUE\_\allowbreak{}KERNEL\_\allowbreak{}POL}, \texttt{OCRT\_\allowbreak{}WATER\_\allowbreak{}MIE\_\allowbreak{}MOMENT\_\allowbreak{}N\_\allowbreak{}MU}, \texttt{OCRT\_\allowbreak{}WATER\_\allowbreak{}VALUE\_\allowbreak{}NPHI}는 raw \texttt{value\_\allowbreak{}phase}를 전달하는 수중 연구 경로에 관한 옵션으로, 고수준 성분 CLI를 자동으로 최신 C 직접 FR631 커널로 바꾸는 스위치가 아니다.

\section{\texttt{campaign\_\allowbreak{}runner}: C 프로세스 캠페인 관리}
\label{sec:auto:06_python_batch:10}
이 도구는 Python RT 엔진과 별개다. 2026-08-22의 마이그레이션 배포 구조 \texttt{runtime-\allowbreak{}root/\allowbreak{}01\_\allowbreak{}OCRT\_\allowbreak{}C/\allowbreak{}}를 전제로 C 실행파일을 실행한다. \textbf{현재 GitHub 저장소 루트를 \texttt{-\allowbreak{}-\allowbreak{}runtime-\allowbreak{}root}로 주어 곧바로 실행하는 도구는 아니다.} \texttt{-\allowbreak{}-\allowbreak{}exe}를 지정해도 작업 폴더는 \texttt{runtime-\allowbreak{}root/\allowbreak{}01\_\allowbreak{}OCRT\_\allowbreak{}C}로 유지된다. 현재 저장소 기반의 새 생산 작업은 \ref{chap:08_input_formats}장과 부록~\ref{chap:04_examples}의 절차를 먼저 따른다.

당시 설정·문서의 “PSSA ON”은 기록의 명칭이다. 실행 명령에는 \texttt{-\allowbreak{}-\allowbreak{}pssa}를 넣는데, \textbf{C v1.11.1에서는 이 스위치가 IPSS를 켠다.} 과거 C 실행파일은 다른 보정을 쓸 수 있으므로 바이너리 버전과 SHA를 반드시 같이 읽는다. 이 도구가 오래된 실행환경을 자동으로 최신 C로 바꿔 주지는 않는다.

\subsection{실행 옵션}
\label{sec:auto:06_python_batch:10.1}

{\TableFont
\begin{longtable}{@{}>{\raggedright\arraybackslash}p{\dimexpr 0.3000\linewidth-5.33pt\relax}>{\raggedright\arraybackslash}p{\dimexpr 0.2300\linewidth-5.33pt\relax}>{\raggedright\arraybackslash}p{\dimexpr 0.4700\linewidth-5.33pt\relax}@{}}
\toprule
\textbf{옵션} & \textbf{기본값} & \textbf{설명} \\
\midrule
\endfirsthead
\toprule
\textbf{옵션} & \textbf{기본값} & \textbf{설명} \\
\midrule
\endhead
\bottomrule
\endfoot
\texttt{-\allowbreak{}h}, \texttt{-\allowbreak{}-\allowbreak{}help} & --- & 도움말 \\
\texttt{-\allowbreak{}-\allowbreak{}runtime-\allowbreak{}root DIR} & 필수 & \texttt{01\_\allowbreak{}OCRT\_\allowbreak{}C} 하위 폴더가 있는 마이그레이션 루트 \\
\texttt{-\allowbreak{}-\allowbreak{}matrix CSV} & 필수 & 실제 C 명령 인자와 검증 조건을 담은 실행 행렬 \\
\texttt{-\allowbreak{}-\allowbreak{}output-\allowbreak{}dir DIR} & 필수 & 결과·로그·매니페스트 저장 폴더 \\
\texttt{-\allowbreak{}-\allowbreak{}config JSON} & 동봉 \texttt{config/\allowbreak{}polarization\_\allowbreak{}campaign\_\allowbreak{}pssa\_\allowbreak{}16core.\allowbreak{}json} & 실행 정책·환경·검증 조건 \\
\texttt{-\allowbreak{}-\allowbreak{}exe PATH} & 자동 탐색 & 실행파일 절대/상대경로. 생략 시 \texttt{01\_\allowbreak{}OCRT\_\allowbreak{}C/\allowbreak{}build/\allowbreak{}}의 \texttt{ocrt\_\allowbreak{}polarization}, \texttt{ocrt}, \texttt{ocrt\_\allowbreak{}v1.\allowbreak{}2}와 \texttt{.\allowbreak{}exe} 변형 탐색 \\
\texttt{-\allowbreak{}-\allowbreak{}workers N} & 설정 파일의 16 & 독립 C 프로세스 수. 현재 정책은 1--16 \\
\texttt{-\allowbreak{}-\allowbreak{}pilot N} & 0, 제한 없음 & 첫 N행만 실행하고 용량·시간 추정 출력 \\
\texttt{-\allowbreak{}-\allowbreak{}max-\allowbreak{}runs N} & 0, 제한 없음 & 실행 행 수 제한. \texttt{-\allowbreak{}-\allowbreak{}pilot} 절단 뒤 적용 \\
\texttt{-\allowbreak{}-\allowbreak{}dry-\allowbreak{}run} & 끄기 & 명령과 실행설정만 생성, C 계산 생략. 바이너리·행렬 검사와 파일 생성은 수행 \\
\texttt{-\allowbreak{}-\allowbreak{}keep-\allowbreak{}raw} & 끄기 & 검증·압축 후에도 원본 CSV 유지 \\
\end{longtable}
}

기본 정책은 worker당 \texttt{OMP\_\allowbreak{}NUM\_\allowbreak{}THREADS=\allowbreak{}1}, \texttt{OCRT\_\allowbreak{}ADVANCED=\allowbreak{}1}, 직접 수중 입자 커널이며 광범위 truncation을 금지한다. 물 실행 행에 \texttt{-\allowbreak{}-\allowbreak{}n-\allowbreak{}mu-\allowbreak{}water}가 없으면 중단한다. \texttt{run\_\allowbreak{}campaign.\allowbreak{}py}는 \texttt{-\allowbreak{}-\allowbreak{}pssa}가 누락되면 추가하지만, 보관 명령을 변환하는 도구는 누락된 명령을 오류로 처리한다. 설정 파일은 정책 기록이며, 기능 지원 여부는 실제로 실행하는 C 버전에서 다시 확인한다.

\subsection{행렬과 실행 예}
\label{sec:auto:06_python_batch:10.2}
행렬의 열은 \texttt{run\_\allowbreak{}id,\allowbreak{}stage,\allowbreak{}uses\_\allowbreak{}water,\allowbreak{}expected\_\allowbreak{}rows,\allowbreak{}output\_\allowbreak{}relpath,\allowbreak{}command\_\allowbreak{}args\_\allowbreak{}json,\allowbreak{}required\_\allowbreak{}columns\_\allowbreak{}json,\allowbreak{}expected\_\allowbreak{}vza\_\allowbreak{}json,\allowbreak{}expected\_\allowbreak{}raa\_\allowbreak{}json,\allowbreak{}notes}이다. \texttt{command\_\allowbreak{}args\_\allowbreak{}json}은 실행파일을 제외한 문자열 인자 목록이며, 셸 전체 명령문이나 파이프가 아니다. 출력 옵션은 실행기가 추가하므로 행렬에 \texttt{-\allowbreak{}-\allowbreak{}output-\allowbreak{}full-\allowbreak{}grid/\allowbreak{}-\allowbreak{}-\allowbreak{}batch-\allowbreak{}full-\allowbreak{}grid/\allowbreak{}-\allowbreak{}-\allowbreak{}output}을 넣지 않는다. \texttt{run\_\allowbreak{}id}는 고유해야 한다. 출력 상대경로는 캠페인 폴더 안의 파일명으로 정한다.

기존 마이그레이션 실행환경이 준비된 경우, \texttt{code/\allowbreak{}campaign\_\allowbreak{}runner}에서:


\begin{ManualCode}
python runner/run_campaign.py --runtime-root runtime/MIGRATION_PKG_2026-08-19 --matrix matrices/PURE_RAYLEIGH_15RUN_MATRIX.csv --output-dir results/rayleigh_reference --workers 4 --pilot 2 --dry-run
python runner/run_campaign.py --runtime-root runtime/MIGRATION_PKG_2026-08-19 --matrix matrices/PURE_RAYLEIGH_15RUN_MATRIX.csv --output-dir results/rayleigh_reference --workers 4 --pilot 2
python runner/monitor_progress.py --output-dir results/rayleigh_reference
\end{ManualCode}

실제 명령·격자·결과를 확인한 뒤 \texttt{-\allowbreak{}-\allowbreak{}pilot 2}를 빼고 같은 출력 폴더로 이어서 실행한다. 압축 결과가 존재하면 구조·행 수·지정 VZA/RAA·선택한 수치 열의 유한성을 다시 검사해 재사용한다. 이 재사용도 모든 입력 설정의 물리적 동일성을 자동 보증하지 않으므로 다른 명령 행렬에는 새 출력 폴더를 쓴다.

\texttt{PURE\_\allowbreak{}RAYLEIGH\_\allowbreak{}15RUN\_\allowbreak{}MATRIX.\allowbreak{}csv}는 5파장 \ensuremath{\times} 3 SZA의 알려진 기준이다. \texttt{HISTORICAL\_\allowbreak{}RUN\_\allowbreak{}MATRIX\_\allowbreak{}REQUIRED.\allowbreak{}csv}는 과거 8,045건을 복원한 완성 행렬이 아니므로 빈 상태의 전체 캠페인을 실행할 수 없다. 누락된 조건을 추측해서 역사적 재현으로 부르지 않는다.

출력은 다음과 같다.


\begin{ManualCode}
results/<campaign>/
  run_config.json             실행환경·입력행렬·바이너리 SHA
  run_commands.txt            실제 C 인자
  run_manifest.csv            실행별 상태·시간·행 수·SHA·로그 경로
  RUN_SUMMARY.json            실행 요약
  PILOT_OR_RUN_ESTIMATE.json  용량·시간 추정
  logs/*.stdout.log
  logs/*.stderr.log
  <stage>/*.csv.gz            검증 후 확정한 원본 결과
\end{ManualCode}

검증을 통과한 CSV를 압축하고 최종 이름으로 바꾼다. 현재 검증기의 유한성 검사는 지정된 일부 출력 열에 한정되므로 모든 열의 과학적 유효성 검사와 동일하지 않다. 시간 추정 파일의 \texttt{estimated\_\allowbreak{}16\_\allowbreak{}worker\_\allowbreak{}wall\_\allowbreak{}s\_\allowbreak{}ideal}은 이름 그대로 16 worker로 나눈 이상적 값이다. \texttt{-\allowbreak{}-\allowbreak{}workers 4} 같은 실제 설정의 예상시간으로 읽지 않는다.

\subsection{보조 도구의 전체 옵션}
\label{sec:auto:06_python_batch:10.3}

{\TableFont
\begin{longtable}{@{}>{\raggedright\arraybackslash}p{\dimexpr 0.3400\linewidth-4.00pt\relax}>{\raggedright\arraybackslash}p{\dimexpr 0.6600\linewidth-4.00pt\relax}@{}}
\toprule
\textbf{도구} & \textbf{옵션} \\
\midrule
\endfirsthead
\toprule
\textbf{도구} & \textbf{옵션} \\
\midrule
\endhead
\bottomrule
\endfoot
\texttt{monitor\_\allowbreak{}progress.\allowbreak{}py} & \texttt{-\allowbreak{}h/\allowbreak{}-\allowbreak{}-\allowbreak{}help}, 필수 \texttt{-\allowbreak{}-\allowbreak{}output-\allowbreak{}dir}. 매니페스트의 상태별 수량·용량·최근 행을 읽음 \\
\texttt{build\_\allowbreak{}matrix\_\allowbreak{}from\_\allowbreak{}run\_\allowbreak{}commands.\allowbreak{}py} & \texttt{-\allowbreak{}h/\allowbreak{}-\allowbreak{}-\allowbreak{}help}, 필수 \texttt{-\allowbreak{}-\allowbreak{}commands}, 필수 \texttt{-\allowbreak{}-\allowbreak{}output}, \texttt{-\allowbreak{}-\allowbreak{}default-\allowbreak{}expected-\allowbreak{}rows} 기본 1296. 보관된 정확한 명령을 행렬과 \texttt{.\allowbreak{}audit.\allowbreak{}csv}로 변환 \\
\texttt{inspect\_\allowbreak{}historical\_\allowbreak{}package.\allowbreak{}py} & \texttt{-\allowbreak{}h/\allowbreak{}-\allowbreak{}-\allowbreak{}help}, 필수 \texttt{-\allowbreak{}-\allowbreak{}archive}, \texttt{-\allowbreak{}-\allowbreak{}output-\allowbreak{}dir} 기본 \texttt{historical\_\allowbreak{}import}. 과거 tar에서 조건·명령 기록과 SHA를 추출 \\
\texttt{setup\_\allowbreak{}runtime.\allowbreak{}py} & \texttt{-\allowbreak{}h/\allowbreak{}-\allowbreak{}-\allowbreak{}help}, 필수 \texttt{-\allowbreak{}-\allowbreak{}parts-\allowbreak{}dir}, 필수 \texttt{-\allowbreak{}-\allowbreak{}runtime-\allowbreak{}dir}, \texttt{-\allowbreak{}-\allowbreak{}verify-\allowbreak{}only}, \texttt{-\allowbreak{}-\allowbreak{}no-\allowbreak{}python-\allowbreak{}data}, \texttt{-\allowbreak{}-\allowbreak{}no-\allowbreak{}build}, \texttt{-\allowbreak{}-\allowbreak{}march} 기본 \texttt{auto} (\texttt{auto/\allowbreak{}cascadelake/\allowbreak{}x86-\allowbreak{}64-\allowbreak{}v3/\allowbreak{}x86-\allowbreak{}64}), \texttt{-\allowbreak{}-\allowbreak{}full-\allowbreak{}mie-\allowbreak{}hash-\allowbreak{}audit}, \texttt{-\allowbreak{}-\allowbreak{}extract-\allowbreak{}workers} 기본 4, \texttt{-\allowbreak{}-\allowbreak{}force} \\
\end{longtable}
}

\texttt{setup\_\allowbreak{}runtime.\allowbreak{}py}는 당시 마이그레이션 기반 파일·정책 overlay·PART00--PART10 묶음을 요구하는 복원 설치기다. 현재 \texttt{data-\allowbreak{}v1} zip 네 개를 받는 \texttt{fetch\_\allowbreak{}data}와 다른 절차다. \texttt{-\allowbreak{}-\allowbreak{}verify-\allowbreak{}only}는 입력 압축파일 검증 후 종료, \texttt{-\allowbreak{}-\allowbreak{}no-\allowbreak{}python-\allowbreak{}data}는 Python 짝 자료 설치 생략, \texttt{-\allowbreak{}-\allowbreak{}no-\allowbreak{}build}는 빌드 생략, \texttt{-\allowbreak{}-\allowbreak{}full-\allowbreak{}mie-\allowbreak{}hash-\allowbreak{}audit}는 추출된 전체 Mie SHA 검사이다. \textbf{\texttt{-\allowbreak{}-\allowbreak{}force}는 이미 존재하는 \texttt{-\allowbreak{}-\allowbreak{}runtime-\allowbreak{}dir}을 삭제하고 다시 만든다.} 원래 보관 자료·결과를 두는 폴더에 쓰지 않는다.

\section{작성·확인 근거}
\label{sec:auto:06_python_batch:11}
\begin{itemize}
\item \href{https://github.com/brtnt/OCRT/blob/54861c68e35ec5aaa3938fad53e93dc2dfa7eb1d/code/OCRT_Python/produce_grid.py}{배치 CLI와 출력 필드}
\item \href{https://github.com/brtnt/OCRT/blob/54861c68e35ec5aaa3938fad53e93dc2dfa7eb1d/code/OCRT_Python/ocrt_solve.py}{단일/순차 격자 CLI}
\item \href{https://github.com/brtnt/OCRT/blob/54861c68e35ec5aaa3938fad53e93dc2dfa7eb1d/code/OCRT_Python/ocrt_fullgrid.py}{각도 LUT CLI}, \href{https://github.com/brtnt/OCRT/blob/54861c68e35ec5aaa3938fad53e93dc2dfa7eb1d/code/OCRT_Python/ocrt_py/lut.py}{각도 LUT 출력}
\item \href{https://github.com/brtnt/OCRT/blob/54861c68e35ec5aaa3938fad53e93dc2dfa7eb1d/code/OCRT_Python/ocrt_py/batch_driver.py}{배치 계산 경로·투과율 정의}, \href{https://github.com/brtnt/OCRT/blob/54861c68e35ec5aaa3938fad53e93dc2dfa7eb1d/code/OCRT_Python/ocrt_py/driver.py}{단일 결합 경로}
\item \href{https://github.com/brtnt/OCRT/blob/54861c68e35ec5aaa3938fad53e93dc2dfa7eb1d/code/OCRT_Python/ocrt_py/constituent.py}{성분 모델·EAP gate}, \href{https://github.com/brtnt/OCRT/blob/54861c68e35ec5aaa3938fad53e93dc2dfa7eb1d/code/OCRT_Python/ocrt_py/backend.py}{배열 backend}
\item \href{https://github.com/brtnt/OCRT/blob/54861c68e35ec5aaa3938fad53e93dc2dfa7eb1d/code/OCRT_Python/gpu_gate.py}{GPU 대조 도구}, \href{https://github.com/brtnt/OCRT/blob/54861c68e35ec5aaa3938fad53e93dc2dfa7eb1d/code/OCRT_Python/verify.py}{과거 C 기준 검사}
\item \href{https://github.com/brtnt/OCRT/blob/54861c68e35ec5aaa3938fad53e93dc2dfa7eb1d/scripts/MIE_SHA256SUMS_data-v1.txt}{현재 배포 자료 매니페스트}
\item \href{https://github.com/brtnt/OCRT/blob/54861c68e35ec5aaa3938fad53e93dc2dfa7eb1d/code/campaign_runner/runner/run_campaign.py}{C 캠페인 실행기}, \href{https://github.com/brtnt/OCRT/blob/54861c68e35ec5aaa3938fad53e93dc2dfa7eb1d/code/campaign_runner/config/polarization_campaign_pssa_16core.json}{기본 설정}
\end{itemize}

이 부록의 검사는 실제 도움말 출력, 입력 파서, 출력 필드, 함수 호출 경로, 로컬 자료 로딩을 대조하고 \ref{sec:auto:06_python_batch:7.1}절의 R2 한 점 실행을 확인한 것이다. GPU 실행·현재 C와의 전체 정밀도 대조·EAP 보조자료가 필요한 R1 실행 성공을 새로 검증했다는 의미는 아니다.

